% Thorne & Campolattaro 1967, ApJ 149, 591
% Page 1 (p. 591) transcription

\title{NON-RADIAL PULSATION OF GENERAL-RELATIVISTIC STELLAR MODELS.
I. ANALYTIC FUNCTIONS FOR $l \geq 2$\,\footnote{* Supported in part by the Office of Naval Research (Nonr-220(47)) and in part by the National Science Foundation (GP-3391).}}

\author{Kip S. Thorne\footnote{\dag\ Alfred P. Sloan Foundation Research Fellow.}\\
\textit{California Institute of Technology, Pasadena}\\
\and
Alfonso Campolattaro\\
\textit{University of California, Irvine}}

\date{Received February 21, 1967}

\maketitle

\begin{abstract}
The theory of small, adiabatic, non-radial perturbations of a star away from hydrostatic equilibrium
is developed within the framework of general relativity. The unperturbed equilibrium configurations are
an arbitrary, non-rotating, general-relativistic stellar model. The departures from equilibrium are
analyzed into tensorial spherical harmonics and then into complex normal modes with various mixtures
of incoming and outgoing gravitational waves. A discussion is given of the expansion of real, physical
pulsations in terms of gravitational radiation in terms of the complex normal modes. Criteria
are developed for stability against non-radial pulsations; and methods are devised for computing numerically the pulsation frequencies, eigenfunctions, and gravitational radiation damping times of the stable,
real quasi-normal modes of pulsation.
\end{abstract}

\section{Motivation}

In the last four years astronomical discoveries and theoretical considerations have
motivated a detailed development of the general-relativistic theory of stellar structure
and dynamics. The discovery of galactic X-ray sources and the development of detailed
hydrodynamic models for supernovae have given impetus to theoretical research on neutron
stars, while supermassive stars have been studied in connection with quasi-stellar radio
sources (QSS's).

One of us (Thorne 1966, 1967) has recently written a detailed review of the general-relativistic
theory of stellar structure and dynamics as it has been developed to date in response to the current
interest in neutron stars and supermassive stars. As was emphasized in that review, one of the most
important next steps in the development of the theory is the analysis of non-radial pulsations of
relativistic stellar models. In this paper we present such an analysis.

Non-radial pulsations are of interest for a number of reasons: If, as current theory suggests,
neutron stars are formed in some supernova explosions by the collapse of the core of a star which is
near the end point of thermonuclear evolution, then, when first formed, a neutron star will pulsate
wildly. The initial energy of pulsation will be of the order of the kinetic energy of collapse, which
will be between $\sim 0.01$ and $\sim 0.2$ of the rest mass-energy of the star.
\cite{Hoyle1964}, \cite{Narikiar1964}, \cite{Cameron1965a}, \cite{Cameron1965b}, and \cite{Finzi1965}
have argued that the gradual transfer of this huge pulsation energy to the supernova envelope which surrounds the collapsing core may have important observational consequences.
(See \cite{Wheeler1966}; \cite{Meltzer1966}; \cite{Thorne1966}; \cite{Finzi1966}; and
\cite{Tsuruta1966} and \cite{Cameron1966} for critiques of and further developments of these ideas.)
In order to evaluate such suggestions and make them more concrete, one needs an understanding of the
theory of both radial and non-radial pulsations of relativistic stellar models.

The non-radial pulsations of neutron stars are of interest also because of their intimate connection
with gravitational radiation. Rough estimates by \cite{Zee1966} and \cite{Wheeler1966}
